\chapter*{Motivation}\label{chap:motivation}

Reinforcement learning (RL) systems learn through interaction: an agent acts in an environment and updates its policy from the rewards it observes. When rewards are dense, the credit assignment problem -- attributing the eventual outcome to the individual actions that produced it -- is largely solved by temporal-difference and policy-gradient methods that propagate signal step by step. The picture changes drastically when feedback is delivered only at the end of an episode. Long horizons, sparse signals, and noisy returns combine into very high gradient variance, slow convergence, and policies that learn from accident more than from understanding.

This work is motivated by a recent line of research that treats the trajectory as a cooperative game, with actions playing the roles of cooperating agents and the terminal reward acting as the value of the grand coalition. The SCAR method \cite{scar2025} demonstrated that distributing the terminal reward across tokens of a generated sequence using Shapley values produces a dense, fair, and theoretically grounded reward decomposition that significantly accelerates RLHF-style fine-tuning of large language models. Crucially, however, SCAR pays a high price for this signal: it relies on a separately trained reward model that approximates human preferences and that must be queried for many counterfactual sequences.

Model-based RL, and the Dreamer family of algorithms \cite{dreamer2020,dreamerv3} in particular, offers a complementary opportunity. Dreamer maintains a learned latent dynamics model and a learned reward predictor, both trained jointly with the policy. These components can be queried offline, in latent space, to evaluate counterfactual continuations of a trajectory \emph{without further interaction with the real environment}. This makes Dreamer a natural host for Shapley-style credit assignment: the world model can supply exactly the kind of cheap counterfactual evaluations that the SCAR estimator demands, and it can do so for arbitrary action sequences rather than only for textual ones.

The aim of this thesis is to bring these two lines of work together. We propose a method that decomposes terminal rewards into per-action contributions using Monte Carlo Shapley sampling, with coalition values estimated through latent imagined rollouts of a Dreamer-style world model. The dense per-action contributions are then fed back as a training signal for both the policy and the dynamics model. We focus on text-based and text--procedural tasks because they exhibit the long, discrete, terminally rewarded structure where credit assignment is hardest, while remaining computationally tractable for prototype-scale experiments.